\section{Prioritized next-stage roadmap}
\label{sec:next-roadmap}

The priorities below are ordered by how directly the existing exact
reductions expose a checkable next theorem.  Priorities one and two
may run in parallel because zero-mode cancellation and energy survival
are both required.

\subsection{Priority 1: zero-mode low-frequency, Tauberian, and
surrogate branch}

\paragraph{Deliverable 1A: literal low-window export.}
Write every
\(\widehat A_X(r)\), \(0<d_{q_X}(r,0)\le R_X\), in the native affine
variables, retain the TPC-86 multiplier, and partition resonant,
short-fiber, and generic sectors without deletion.  The first theorem
target is the complete multiplier-weighted window, not an individual
generic frequency.  Choose \(R_X=X^{\varrho+o(1)}\) and filter order
\(k\) so that \(R_X<q_X/2\),
\(\varrho<267/400\), and
\((k-\tfrac12)\varrho>133/400+\eta_Z\); the
\(133/600+\eps\) window, with fixed \(0<\eps<107/240\), is
reserved for the second-order zero-loss corner unless
\(\eta_Z<3\eps/2\).

\paragraph{Deliverable 1B: signed Tauberian certificate.}
For any logarithmic input, compute the exact unnormalized primitive
and its signed backward Ces\`aro defect on every physical prefix.
Then prove that the coefficient-weighted fiber defects, short-fiber
return, and content completion together meet
\eqref{eq:tauberian-aggregate-target}, or the corresponding
quantitative target of the attached block.  A per-fiber \(o(1)\)
defect is not enough when the number or weights of fibers grow.  If
only normalized window cancellation is available, record a negative
transfer result rather than treating it as ordinary cancellation.

\paragraph{Deliverable 1C: surrogate feasibility test.}
Construct candidate same-frame surrogates only after an exact support
map.  Prove or disprove the simultaneous
\[
 \rho(b_X),\quad
 \|A_X-b_X\|_2/\|b_X\|_2
 \le X^{o(1)}/J_X
\]
requirements.  A sharp lower bound on the best achievable error would
be a publishable negative result.

\subsection{Priority 2: determinant mass and equal-determinant
coherence}

\paragraph{Deliverable 2A: native-key census.}
Implement the TPC-84 streaming/external-sort schema with exact
determinant and provenance keys.  Report \(M_X,D_X,E_{\rm trip}\),
\(C_{\rm eq}\), dominant-fiber mass, pairing-class totals, and the
unpaired remainder under one normalization.  This computation chooses
the next theorem; it does not prove asymptotics.

\paragraph{Deliverable 2B: one growing lower bound.}
Target either:
\[
 M_X\ge X^{-o(1)}N_{0,X},
\]
a positive same-scale lower bound for one literal pairing class with
controlled remainder, or a rigorously quantified fixed-loss profile.
If a fixed-power mass loss is proved instead, abandon the natural
corner and recalculate the full TPC-88 ledger.

\subsection{Priority 3: primitive affine arithmetic}

TPC-90 shows that the next useful theorem is not another local root
count.  Expand one literal low-frequency or determinant block into
the squarefree-supported Liouville correlation and seek a bound that
is uniform in the growing slopes, intercepts, origins, masks, and
prefixes.  The preferred progression is:
\begin{enumerate}[label=(\roman*),leftmargin=2.2em]
\item one literal residual block with a fixed-power or arbitrary-log
  saving;
\item a controlled family of square-divisor moduli with exact
  complementary return;
\item the full multiplier-weighted low window or TPC-18 generic
  block; and
\item complete outer reassembly.
\end{enumerate}
At each stage, averaged shifts, frozen coefficients, and density-one
ranges retain those qualifiers.

\subsection{Priority 4: full-block reassembly, exact cover first}

Export the TPC-16/TPC-18 hard packet to the TPC-91 incidence format.
The first calculation is exact feasibility, preferably with rational
or interval-certified arithmetic.  There are three publishable
outcomes:
\begin{enumerate}[label=(\alph*),leftmargin=2.2em]
\item an allowed exact cover with acceptable coefficient cost;
\item a dual residual or Farkas certificate that kills the proposed
  zero-residual block family; or
\item an explicit residual plus a new physical theorem controlling
  it.
\end{enumerate}
Outcome (a) closes reassembly with zero residual.  Outcome (b) does
not kill the original TPC-18 residual route.  Outcome (c) closes
reassembly only at the rate actually proved.  Primitive determinant
dispersion remains a separate subsequent theorem.

\subsection{Final merge audit}

No branch is declared complete in isolation.  Before any endpoint
claim, rerun the following audit:
\begin{enumerate}[label=(\arabic*),leftmargin=2.2em]
\item literal physical coefficients and one prescribed \(h_0\);
\item actual support and all native keys;
\item one global atomic normalization;
\item canonical, complete representation with no deleted sectors;
\item the recorded determinant compatibility
  \(\lambda_D\le2\eta_Z\); and
\item the independent strict physical budget
  \(\Lambda_{\rm phys}<1/400\).
\end{enumerate}
Then verify the TPC-18 reassembly (exact cover or a controlled
physical residual) and primitive determinant input separately.  Any
duplicated loss is removed before the comparison.

\begin{proposition}[Roadmap outcome trichotomy]
\label{prop:roadmap-trichotomy}
Every prioritized branch has three honest outcomes:
\begin{enumerate}[label=(\roman*),leftmargin=2.2em]
\item a literal positive estimate, which advances the dependency
  graph at its proved level;
\item an exact negative certificate, which permanently removes the
  specified route or family; or
\item a conditional interface or numerical diagnostic, which is
  retained without promotion.
\end{enumerate}
None of these labels implies a prime-pair theorem unless the entire
remaining conjunction is proved.
\end{proposition}

\begin{proof}
The positive and negative certificates are the go/stop tests in
\cref{sec:stop-go-tests}; every other outcome lacks at least one
promotion item in \cref{prop:promotion-rule}.  The final statement
follows because the completion conjunction remains unproved.
\end{proof}
